\section{Artifact and Discussion}
\label{sec:artifact}

\subsection{Artifact}

The executable artifact separates the theorem-facing model, empirical enlargement, benchmark runner, imported references, and Lean finite bridge:
the core paths are \texttt{mino/models/}, \texttt{mino/data/}, \texttt{scripts/}, \texttt{lean/MiNO/Microlocal/}, and \texttt{manuscript/jmlr/}.
The manuscript-source bundle and executable artifact may be distributed separately; an artifact-inclusive bundle contains these paths together with raw JSON/CSV summaries, seed manifests, and exact runner commands.  The Lean component verifies the finite packet landing layer once the classical microlocal estimates have supplied a packet matrix with tails and budgets.  The Python artifact writes raw JSON, aggregate CSV, and manifest files so that benchmark rows can be audited and resumed.

\subsection{Closed Results}

The finite landing spine is closed for the stated packet model: transported-symbol perturbation, bounded-stencil stability, shell/tube truncation, cross-resolution transfer, and the finite landing layer are checked or replayed in the Lean companion.  The operator-level theorem spine is the $\ell^2$ sparse packet-matrix transfer coupled to the continuous no-caustic FIO reduction: short-time half-wave propagation gives a canonical graph, non-stationary phase gives off-tube decay, and packet sparsity converts that decay into a bounded-stencil approximation.  Together these pieces justify carrier-bound \MiNO-Core as a theorem-guided architecture for short-time no-caustic wave propagation.

At the executable level, the theorem-facing \texttt{spectral\_order} branch is aligned with the finite symbol statement: it is a pure multiplier on packet coefficients, while additive feature shifts are assigned to the empirical branch or to residual/landing terms.  For high-$k$ Helmholtz, the artifact exposes \texttt{helmholtz\_resolvent}, a near-characteristic radiation proxy corresponding to Corollary~\ref{cor:helmholtz-resolvent-symbol-budget}.  This branch uses an auto-calibrated shell, signed real/imaginary resolvent weights, an outgoing gate, and a bias-free latent complex-pair action for the signed multiplier.  Repository tests check zero-input equivariance of the theorem-facing symbol output, exact finite real multiplication by complex pairs, and differentiability of the symbol-order proxy with respect to packet values.

\subsection{Mechanism Evidence}

The carrier-bound \texttt{branch\_id\_v3} and \texttt{branch\_id\_v3\_controls} campaigns supply the central mechanism evidence.  They keep final metadata batch-specific, apply transported synthesis by a packet-energy-weighted warp, transport a high-frequency input carrier from the original packet coefficients, low-pass restrict skip and refinement paths, and disable token-level refinement in the binding profile.  Removing metadata-flow supervision leaves field error essentially unchanged; removing transport consistently worsens field error; randomizing metadata worsens it further.  This is the empirical alignment point between the sparse packet-matrix theorem and the carrier-bound executable model.

\subsection{Evidence Boundary}

The completed evidence has a clear center: the sparse canonical packet matrix is a field-level mechanism on controlled short-time wave probes.  The theory identifies the packet operator class, the architecture implements a carrier-bound realization, and the benchmark corrupts the corresponding branch.  The conclusion is primitive-level: when high-frequency content must be landed after nontrivial canonical motion, explicit packet transport is a directly testable representation of the law.  Same-refinement UNO/WNO rows show that strong multiscale baselines can learn useful transported behavior implicitly; \MiNO's contribution is to expose and test that law directly.

The high-$k$ Helmholtz and Navier--Stokes rows serve as calculus stress paths.  For Helmholtz, the local outgoing resolvent requires anisotropic packet resolution, branch separation, phase/landing accuracy, and a near-characteristic radiation symbol.  The \texttt{helmholtz\_branched\_highk} probe shows that branch multiplicity alone is not the active budget at the current scale; the \texttt{helmholtz\_highk\_careful} profile points instead to anisotropic carrier-bound resolvent reconstruction.  For Navier--Stokes-style dynamics, the relevant extension is an advective FIO branch coupled to a dissipative $\Psi$DO branch and pressure/closure residuals.  These rows motivate the extension ledger in Section~\ref{sec:admissible-extension-axes}.

High-$k$ Helmholtz is the most direct future flagship benchmark for the microlocal primitive, because it stresses phase coherence, outgoing radiation, and frequency-scaling rather than generic smooth interpolation.  A solver-level claim requires the stronger contract in Section~\ref{sec:branched-helmholtz}: increasing-$k$ rows, phase and radiation residuals, reference-solver accuracy, and comparisons against both general neural operators and specialized Helmholtz solvers.  Navier--Stokes is a secondary nonlinear stress path; its microlocal hook is transport--diffusion splitting, not a standalone turbulence theory.

Three implementation axes are separated from the main wave theorem.  The separable St\"ormer--Verlet metadata branch gives the cleanest symplectic metadata statement; the main theorem itself only requires a learned near-canonical map with the displayed compact-window error budget.  The local symbol branch is secondary on transport-dominant wave probes and is naturally tested in identity-canonical elliptic, dissipative/parabolic, and outgoing-resolvent regimes.  The learned landing decoder is treated as an admissible finite implementation layer for $E_{\mathrm{land},\theta}$, $E_{\mathrm{tsyn}}$, and $E_{\mathrm{car}}$; transport corruption remains visible after the decoder is enabled.

The extension ledger records how additional microlocal structure enters the same packet budget.  Anisotropic packets, finite multibranch canonical relations, stronger $S^m$ symbol control, outgoing resolvent-symbol enforcement, and nonlinear transport--diffusion coupling each has a packet-calculus object and an executable hook.  Once the hook has a verified finite defect term, it enters the packet-budget theorem additively.

Two foundation directions are separated from the present paper.  An Egorov-style neural calculus would test whether learned operators conjugate pseudodifferential probes by the learned canonical map; the finite objects are $\rho_{\mathrm{Eg},h}^{\mathrm{pkt}}$ for a frozen linear Core and $\rho_{\mathrm{Eg},h}^{\mathrm{Jac}}(u,j)$ for the nonlinear executable model.  An operator-microscopy program would recover and compile a microlocal fingerprint, in the no-caustic graph case $\mathfrak F(T)=(\kappa,S,a,R)$, consisting of canonical map, action/phase primitive, amplitude, and residual.  Multi-chart caustics and Maslov transition data require a separate Lagrangian atlas.

\subsection{Conclusion}

\MiNO\ advances a testable microlocal design principle: neural operators for oscillatory PDEs should expose phase-space packet structure rather than rely only on point-space correlation or global Fourier filtering.  The mathematical contribution is a sparse canonical packet-matrix primitive with a finite machine-checkable landing layer and a short-time wave/FIO approximation theorem for carrier-bound \MiNO-Core.  The empirical contribution is a mechanism-identifiability result: in the carrier-bound architecture, high-frequency reconstruction is tied to transported packet metadata, and transport corruption consistently degrades controlled wave rows.

The resulting claim is specific.  The theorem identifies the operator class, the architecture implements one carrier-bound realization of the transported packet primitive in field synthesis, and the benchmark contract tests whether the trained model uses that primitive.  The high-frequency Helmholtz and Navier--Stokes rows define separate calculus hooks---branch/resolvent structure for outgoing oscillatory inverses and transport--diffusion splitting for nonlinear dynamics---rather than broadening the main mechanism claim.
